\section{Задание 1}

\subsection{Методика измерения}

Для оценки размеров L1d, L2 и L3 программой \texttt{cache\_levels.cpp} была измерена задержка перехода по указателю при разных размерах рабочего набора. L1d хранит данные, а L2 и L3 хранят данные и инструкции. При промахе в L1d данные ищутся на более медленных уровнях, вплоть до оперативной памяти. После прогрева набор, помещающийся в кеш, обычно обслуживается этим или более быстрым уровнем. При превышении ёмкости растёт доля обращений к следующему уровню и увеличивается задержка.

Узлы размером и выравниванием по 64 байта были размещены по одному на строку кеша и связаны в цикл в случайном порядке, чтобы затруднить аппаратную предвыборку (листинг~\ref{lst:cache-levels-cycle}). Каждый следующий адрес определяется чтением указателя из текущего узла. Поэтому измерялась задержка зависимых загрузок, а не пропускная способность при параллельных загрузках (листинг~\ref{lst:cache-levels-chase}).

\verbatimlistingcaption{Структура узла и построение случайного цикла}{lst:cache-levels-cycle}
\begin{Verbatim}[frame=single,numbers=left,fontsize=\small,breaklines=true]
Узел (64 байта, выравнивание по 64 байтам):
    следующий: указатель на узел

Алгоритм: построить_цикл(узлы)
Вход: непустой массив узлов.
Выход: адрес первого узла построенного цикла.

n <- число узлов
порядок <- перемешать индексы от 0 до n - 1
для i от 0 до n - 1:
    текущий <- порядок[i]
    следующий <- порядок[(i + 1) по модулю n]
    узлы[текущий].следующий <- адрес(узлы[следующий])
вернуть адрес(узлы[порядок[0]])
\end{Verbatim}

\verbatimlistingcaption{Обход цепочки указателей}{lst:cache-levels-chase}
\begin{Verbatim}[frame=single,numbers=left,fontsize=\small,breaklines=true]
Алгоритм: обойти(текущий, шаги)
Вход: адрес начального узла и число переходов.
Выход: адрес узла после заданного числа переходов.

повторить шаги раз:
    текущий <- текущий.следующий
вернуть текущий
\end{Verbatim}

В случайном порядке были проверены степени двойки от 1~КиБ до 128~МиБ и промежуточные размеры $1{,}25S$ и $1{,}5S$ между соседними степенями $S$ и $2S$. Для каждого размера было проведено 15 испытаний. Перед каждым измерением были выполнены четыре полных прогревочных обхода. Число переходов $N$ было округлено вверх до ближайшего кратного длине цикла, не меньшего 5~млн. Задержка за шаг рассчитывалась как $(t_1-t_0)/N$, где $t_0$ и $t_1$ были сняты с монотонных часов непосредственно до и после обхода. Числитель выражен в наносекундах. Во всех трёх программах итоговый адрес записывался в переменную \texttt{volatile}, чтобы обход не был удалён при оптимизации. На графике показаны медианы по 15 испытаниям (листинг~\ref{lst:cache-levels-measure}).

\verbatimlistingcaption{Измерение задержки для разных размеров рабочего набора}{lst:cache-levels-measure}
\begin{Verbatim}[frame=single,numbers=left,fontsize=\small,breaklines=true]
Алгоритм: эксперимент с размерами рабочего набора
Вход: заданные в коде границы размеров и параметры испытаний.
Выход: размеры, результаты испытаний и статистики в нс/шаг.

размеры <- пустой список
для S = 1 КиБ, 2 КиБ, 4 КиБ, ..., 128 МиБ:
    добавить S в размеры
    если S < 128 МиБ:
        добавить S + S / 4 и S + S / 2 в размеры

результаты <- пустой список
для каждого размера из перемешанного списка размеров:
    n <- размер в байтах / 64
    узлы <- массив из n узлов
    начало <- построить_цикл(узлы)
    шаги <- n * округлить_вверх(5000000 / n)
    задержки <- пустой список
    повторить 15 раз:
        повторить 4 раза для прогрева:
            адрес <- обойти(начало, n)
            сохранить адрес
        t0 <- показание монотонных часов
        адрес <- обойти(начало, шаги)
        t1 <- показание монотонных часов
        сохранить адрес
        задержка <- (t1 - t0) в наносекундах / шаги
        добавить задержку в задержки
    отсортированные <- отсортировать копию задержек
    медиана <- восьмое значение отсортированного списка
    минимум, максимум <- первое и последнее значения списка
    добавить размер, задержки и статистики в результаты

вывести результаты по возрастанию размера
\end{Verbatim}

\subsection{Результаты измерений}

На рисунке~\ref{fig:cache-levels} представлена зависимость медианной задержки от размера рабочего набора в логарифмическом масштабе по обеим осям.

\begin{figure}[htbp]
    \centering
    \begin{tikzpicture}
        \begin{loglogaxis}[
                width=0.98\linewidth, height=7cm,
                xlabel={Размер рабочего набора, КиБ},
                ylabel={Задержка, нс/шаг},
                log basis x=2,
                xtick={1,8,32,128,512,4096,16384,131072},
                xticklabels={1,8,32,128,512,4096,16384,131072},
                ytick={1,2,5,10,20,50,100},
                yticklabels={1,2,5,10,20,50,100},
                tick label style={font=\small},
                grid=major,
            ]
            % Медианы замеров: размер в КиБ, задержка в нс/шаг.
            \addplot[blue, thick, mark=*, mark size=1pt] coordinates {
                    (1,0.95) (1.25,0.95) (1.5,0.95)
                    (2,0.95) (2.5,0.95) (3,0.95)
                    (4,0.96) (5,0.94) (6,0.96)
                    (8,0.95) (10,0.95) (12,0.95)
                    (16,0.96) (20,1.00) (24,1.06)
                    (32,0.95) (40,1.77) (48,2.12)
                    (64,2.54) (80,2.69) (96,2.58)
                    (128,2.70) (160,2.74) (192,2.93)
                    (256,2.89) (320,3.35) (384,3.51)
                    (512,4.63) (640,6.52) (768,10.15)
                    (1024,10.39) (1280,10.00) (1536,9.75)
                    (2048,13.19) (2560,12.61) (3072,15.50)
                    (4096,35.45) (5120,56.51) (6144,65.06)
                    (8192,74.26) (10240,79.89) (12288,82.84)
                    (16384,89.77) (20480,92.96) (24576,93.03)
                    (32768,92.78) (40960,97.73) (49152,96.34)
                    (65536,97.90) (81920,99.09) (98304,105.69)
                    (131072,101.23)
                };
        \end{loglogaxis}
    \end{tikzpicture}
    \caption{Медианная задержка доступа к данным}
    \label{fig:cache-levels}
\end{figure}

Размеры кешей были оценены по трём областям роста задержки на графике.

\begin{itemize}
    \item Для L1d при 32~КиБ получена медиана $0{,}95$~нс/шаг, характерная для малых наборов, а при 40~КиБ --- уже $1{,}77$~нс/шаг. В качестве оценки принято 32~КиБ, последнее проверенное значение до скачка.
    \item Для L2 при увеличении набора с 384 до 512~КиБ зафиксирован рост медианы с $3{,}51$ до $4{,}63$~нс/шаг, а при 640~КиБ получено $6{,}52$~нс/шаг. За ориентир принято 512~КиБ.
    \item Для L3 при увеличении набора с 3 до 4~МиБ зафиксирован рост медианы с $15{,}50$ до $35{,}45$~нс/шаг, а при 5~МиБ получено $56{,}51$~нс/шаг. За ориентир принято 4~МиБ, где наблюдается резкий рост задержки.
\end{itemize}

Ориентировочные оценки сопоставлены в таблице~\ref{tab:cache-levels} с размерами отдельных экземпляров кеша из столбца \texttt{ONE-SIZE} вывода \texttt{lscpu -C}.

\begin{table}[H]
    \centering
    \caption{Сравнение оценок размеров кеша}
    \label{tab:cache-levels}
    \begin{tabular}{|l|c|c|}
        \hline
        Кеш & Оценка по графику & Размер по \texttt{lscpu} \\
        \hline
        L1d & $\approx 32$~КиБ  & 32~КиБ                   \\
        L2  & $\approx 512$~КиБ & 512~КиБ                  \\
        L3  & $\approx 4$~МиБ   & 4~МиБ                    \\
        \hline
    \end{tabular}
\end{table}
